\scanpage{435}
\mspage{347}

Soit donc $\pi'>\pi$ un autre cardinal, quand peut-on assurer que \inserted{si} $\underline S$ est $\unsure{\Phi|\mathscr I|\mathscr D}$-$\pi$-accessible dans $\unsure{M|\mathcal M}$, il soit aussi $\Phi$-$\pi'$-accessible ?

Les conditions 2) 3) 4) sont visiblement stables par passage de $\pi$ à $\pi'$. Restent les conditions 1) et 5). On sait que si ${\pi'}^{\pi}=\pi'$, alors la condition 1) reste valable pour $\pi'$. Plus \unsure{généralement|généralemt}, si $\pi_0^{\deleted{\unsure{\prime|*}}}$ est un cardinal $\leq\pi$ tel que \unsure{lui|le} soit utile pour $M$ et pour $\underline S$, alors si $\pi'>\pi_0$ est tel que ${\pi'}^{\pi_0}=\pi'$, alors $\pi'$ est aussi utile pour $M$ et pour $\underline S$.

Si on admet l'\unsure{HCC|HCG} (hypothèse du continu généralisée), il en est ainsi pour tt $\pi'>\pi_0$, donc pour \unsure{tt|tout} $\pi'>\pi$.

Reste la condition 5), pour laquelle une réponse est donnée par 4.12.1 (p. 27\deleted{\unsure{2|3}}\unsure{3|2}). Si $\pi_0$ \inserted{$\leq\pi$} est choisi de façon que $f$ soit $\pi_0$-accessible (\unsure{p.e.|p.ex.} $\pi_0=\pi$), et si \deleted{$\pi_1\geq\pi$ $(\unsure{>|\geq}\pi_0)$ \unsure{on pose} est tel que $\pi_1=\sup_{\unsure{a}} f(\unsure{\in|{}<}M_{\pi_0})$ \unsure{ou} $\pi_1=\mathrm{Card}\,\unsure{\mathrm{red}|\mathrm{rel}|\mathrm{réel}}\,\underline S_\pi$,}
\[
\pi'_1=\mathrm{Card}\,\unsure{\mathrm{red}|\mathrm{rel}|\mathrm{réel}}\,\underline S_{\pi_0}\quad(\leq\mathrm{Card}\,\mathrm{red}\,\underline S(\pi)),
\]
\deleted{on voit que}
\[
\pi''_1=\mathrm{Sup}\,(\pi'_1,\pi),
\]
on voit que si $\pi'\geq\pi''_1$ (i.e. $\pi'\geq\pi'_1,\pi$) est tel que $\deleted{\unsure{\pi'\geq\lambda|\pi>\lambda|\pi'}}\ {\pi'}^{\pi_0}=\pi'$, alors on a bien
\[
 f(\underline S_{\pi'})\subset M_{\pi'}.
\]

% archive sheet number: 434

